\chapter{Zusammenfassung und Ausblick}
\label{kap:zusammenfassung}

\section{Zusammenfassung}

Diese Arbeit hat die Verteilung von Arbeitsgängen auf Mitarbeitende bei
Erpolino als Optimierungsproblem formuliert und zwei Wege verglichen,
es zu lösen.

Das Problem ließ sich als ganzzahliges lineares Programm formulieren, in
der Notation des Scheduling als $R\,|\,M_j\,|\,C_{\max}$. Für das
betrachtete Fertigungslos aus \ZahlAufgaben{} Arbeitsgängen und
\ZahlMitarbeiter{} Mitarbeitenden liefert \ac{glpk} in \MilpSekunden{}
Sekunden eine beweisbar optimale Gesamtdauer von \Optimum{} Minuten.
Gegenüber der bisherigen Praxis, die ungefähr einer Greedy-Regel
entspricht und \MakespanGreedy{} Minuten benötigt, ist das eine
Ersparnis von rund einer halben Stunde je Los.

Die Forschungsfrage lässt sich damit beantworten: Ja, die Verteilung
lässt sich sinnvoll optimieren, und der Ertrag ist messbar. Der
überraschende Teil der Antwort betrifft das Wie.

\section{Stärken und Schwächen der Verfahren}

\subsection{Das exakte Verfahren}

Seine Stärke ist die \Index{Gütegarantie}. Der Löser liefert nicht nur
eine Lösung, sondern den Beweis, dass es keine bessere gibt. Diese
Aussage kann keine Heuristik machen, und sie ist in der Diskussion mit
der Arbeitsvorbereitung viel wert: Sie beendet die Frage, ob nicht doch
noch etwas herauszuholen wäre.

Hinzu kommt die Beschreibbarkeit. Das Modell in
Anhang~\ref{anh:mathprog} umfasst rund 20 Zeilen und ist ohne
Programmierkenntnisse lesbar. Ändert sich eine Anforderung, ändert sich
eine Nebenbedingung -- nicht ein Algorithmus.

Seine Schwäche ist das Laufzeitverhalten. Es ist im schlechtesten Fall
exponentiell, und Abschnitt~\ref{sec:skalierung} zeigt, dass der
schlechteste Fall nicht theoretisch bleibt: Ab 60~Arbeitsgängen erreicht
\ac{glpk} die Zeitschranke, ohne die Optimalität zu belegen. Zudem ist
die Rechenzeit nicht vorhersagbar -- die Instanz mit 40~Arbeitsgängen
war leichter als die mit 24.

\subsection{Simulated Annealing}

Seine Stärke ist die planbare Rechenzeit. Sie hängt an der Zahl der
Iterationen, nicht an der Instanz, und wächst über den untersuchten
Bereich nur schwach. Wo der Löser aufgibt, liefert das Verfahren
weiterhin brauchbare Lösungen, bei 80~Arbeitsgängen sogar bessere als
der Löser in 30~Sekunden.

Seine Schwächen sind gravierender, als es die Zahlen nahelegen. Es gibt
keine Gütegarantie: Ohne Vergleichswert bleibt unbekannt, wie weit eine
Lösung vom Optimum entfernt ist. Genau dieser Vergleichswert fehlt aber
in dem Bereich, in dem das Verfahren gebraucht wird -- dort, wo der
Löser versagt. Das ist kein Detail, sondern die eigentliche Schwäche der
Methode.

Hinzu kommt die Empfindlichkeit gegenüber Entscheidungen, die nicht aus
dem Problem folgen. Abschnitt~\ref{sec:energiewirkung} hat gezeigt, dass
die Wahl der Energie über Erfolg und Misserfolg entscheidet: Mit der
naheliegenden rohen Zielgröße erreicht kein einziger Lauf das Optimum,
mit dem Glättungsterm gelingt es. Abkühlplan, Nachbarschaft und
Iterationszahl wirken ähnlich. Nichts davon steht im Modell; alles davon
muss erprobt werden.

\subsection{Empfehlung}

Für Erpolino in der heutigen Größenordnung lautet die Empfehlung
eindeutig: das exakte Verfahren. Es ist schneller, besser und liefert
eine Garantie. Die Heuristik ist die aufwendigere Lösung des leichteren
Problems.

Das ist ein negatives Ergebnis für den Hauptteil dieser Arbeit, und es
ist keines, das sich verstecken lässt. Es ist zugleich das nützlichste
Ergebnis: Ein Unternehmen, das die Zuordnung verbessern will, sollte
nicht mit einer Metaheuristik beginnen. Es sollte mit einem Modell
beginnen.

Die Heuristik wird interessant, sobald die Instanz wächst -- durch mehr
Arbeitsgänge, weitere Standorte oder einen längeren Planungshorizont.
Der gemessene Umschlag bei etwa 60~Arbeitsgängen gilt für \ac{glpk}; mit
einem kommerziellen Löser dürfte er deutlich später liegen.

\section{Ausblick}

Drei Richtungen bieten sich an.

\emph{Das Modell erweitern.} Die Annahmen aus
Abschnitt~\ref{sec:annahmen} sind spürbare Vereinfachungen.
\Index{Rüstzeiten}, Reihenfolgeabhängigkeiten und Fälligkeitstermine
fehlen. Termine ließen sich als Nebenbedingung ergänzen, ohne die
Struktur zu ändern; Reihenfolgeabhängigkeiten führen dagegen auf ein
anderes Problem -- ein Job-Shop --, das erheblich schwerer ist
\parencite{pinedo2016scheduling}.

\emph{Die Zielgröße hinterfragen.} Der Makespan behandelt alle Aufgaben
gleich. In der Praxis sind manche Aufträge dringender als andere. Eine
gewichtete Zielgröße oder die Summe der Verspätungen wäre näher an dem,
was die Arbeitsvorbereitung tatsächlich abwägt.

\emph{Die Datengrundlage verbreitern.} Die Effizienzwerte der
Qualifikationsmatrix stammen aus Erfahrungswerten. Ob sie stimmen, ist
offen, und das Ergebnis hängt unmittelbar an ihnen. Aus dem
\ac{erp}-System ließen sich Ist-Zeiten gewinnen und die Matrix
empirisch schätzen. Ohne diesen Schritt optimiert das Verfahren
zuverlässig auf Zahlen, deren Herkunft niemand belegen kann -- und das
ist am Ende das größere Risiko als die Wahl des Verfahrens.
